\section{Token types}
\label{app:tags}

\begin{table}[!htb]
\centering
\footnotesize
\begin{tabular}{@{}lp{0.72\textwidth}@{}}
\toprule
Token type & A target token has this type when \\
\midrule
Exact repeat, far & the longest earlier match of its preceding context has at least 8 tokens, the
token after the most recent occurrence of that match is the target, and that occurrence lies more
than 512 tokens back \\
Exact repeat, near & the same, with the occurrence at most 512 tokens back \\
Consistent cue & the longest match has 2 to 7 tokens and every earlier continuation of it equals
the target \\
Latest value & the longest match has 2 to 7 tokens and at least two distinct earlier
continuations, and the most recent one is the target \\
Older value & as for latest value, but the target is an earlier continuation other than the most
recent one \\
Novel token & the target token has not appeared earlier in the window \\
No match & not even the one token before the target has an earlier occurrence \\
Near-duplicate line, far & there is no exact match of 8 tokens, the current line has at least 4
distinct token bigrams, and it shares at least half of them (Jaccard) with a line that ends at
least 512 tokens before it and contains the target \\
Closing bracket, far & the target closes a bracket that was opened more than 64 tokens back \\
Next integer & the target is an integer $n$ and the nearest earlier integer within 200 tokens is
$n - 1$ \\
Format & there is no exact match of 8 tokens, the target is punctuation or whitespace, and the
most recent earlier position preceded by the same sequence of token classes (the 5 to 12 tokens
before the target in its line) continued with the same token \\
\bottomrule
\end{tabular}
\caption{Token types used in Figure~\ref{fig:types} and Tables~\ref{tab:a3} and~\ref{tab:types}. Types are computed from the
tokens of each evaluation window alone and can overlap; their shares therefore do not sum to 100\%.
Match lengths are measured on the ladder 1, 2, 3, 4, 6, 8, 12, 16, 24, 32, 48 and 64, and a match
of 2 to 7 tokens is therefore evaluated at length 2, 3, 4 or 6. We
fixed the definitions before the dose grid ran.}
\label{tab:tags}
\end{table}

\begin{table}[!htb]
\centering
\footnotesize
\setlength{\tabcolsep}{4pt}
% Generated by scripts/paper_results.py from runs/. Do not edit.
\begin{tabular}{@{}lrccc@{}}
\toprule
& & GDN + \ingram{} & GDN + \ingram{} & 7:1 hybrid + \ingram{} \\
Token type & \% & 164M tokens & 328M tokens & 1.31B tokens \\
\midrule
All tokens & 100.0 & $-0.142$ & $-0.072$ & $-0.063$ \\
 &  & {\scriptsize [$-0.187$, $-0.102$]} & {\scriptsize [$-0.102$, $-0.046$]} & {\scriptsize [$-0.084$, $-0.044$]} \\
\midrule
Exact repeat, far & 16.4 & $-0.515$ & $-0.337$ & $-0.216$ \\
 &  & {\scriptsize [$-0.642$, $-0.401$]} & {\scriptsize [$-0.429$, $-0.252$]} & {\scriptsize [$-0.290$, $-0.152$]} \\
Exact repeat, near & 45.2 & $-0.045$ & $-0.016$ & $-0.009$ \\
 &  & {\scriptsize [$-0.070$, $-0.027$]} & {\scriptsize [$-0.026$, $-0.009$]} & {\scriptsize [$-0.014$, $-0.005$]} \\
Consistent cue & 12.1 & $-0.202$ & $-0.111$ & $-0.064$ \\
 &  & {\scriptsize [$-0.261$, $-0.156$]} & {\scriptsize [$-0.151$, $-0.078$]} & {\scriptsize [$-0.090$, $-0.044$]} \\
Latest value & 2.3 & $-0.104$ & $-0.037$ & $-0.038$ \\
 &  & {\scriptsize [$-0.133$, $-0.075$]} & {\scriptsize [$-0.059$, $-0.017$]} & {\scriptsize [$-0.049$, $-0.027$]} \\
Older value & 2.1 & $+0.018$ & $+0.089$ & $-0.013$ \\
 &  & {\scriptsize [$-0.026$, $+0.060$]} & {\scriptsize [$+0.042$, $+0.138$]} & {\scriptsize [$-0.027$, $+0.000$]} \\
Novel token & 5.4 & $-0.216$ & $-0.163$ & $-0.112$ \\
 &  & {\scriptsize [$-0.283$, $-0.173$]} & {\scriptsize [$-0.239$, $-0.117$]} & {\scriptsize [$-0.155$, $-0.079$]} \\
No match & 5.4 & $-0.110$ & $-0.063$ & $-0.063$ \\
 &  & {\scriptsize [$-0.166$, $-0.073$]} & {\scriptsize [$-0.117$, $-0.029$]} & {\scriptsize [$-0.088$, $-0.043$]} \\
Near-duplicate line, far & 3.7 & $-0.117$ & $-0.040$ & $-0.051$ \\
 &  & {\scriptsize [$-0.196$, $-0.060$]} & {\scriptsize [$-0.082$, $-0.007$]} & {\scriptsize [$-0.088$, $-0.023$]} \\
Closing bracket, far & 0.1 & $-0.199$ & $-0.116$ & $-0.178$ \\
 &  & {\scriptsize [$-0.448$, $-0.001$]} & {\scriptsize [$-0.302$, $+0.029$]} & {\scriptsize [$-0.295$, $-0.095$]} \\
Next integer & 1.0 & $-0.094$ & $-0.057$ & $-0.055$ \\
 &  & {\scriptsize [$-0.188$, $-0.028$]} & {\scriptsize [$-0.112$, $-0.018$]} & {\scriptsize [$-0.126$, $-0.010$]} \\
Format & 4.7 & $+0.016$ & $+0.012$ & $-0.007$ \\
 &  & {\scriptsize [$+0.004$, $+0.038$]} & {\scriptsize [$+0.005$, $+0.020$]} & {\scriptsize [$-0.013$, $-0.004$]} \\
\bottomrule
\end{tabular}

\caption{Loss of the model with \ingram{} minus loss of the 3:1 hybrid, by token type, for
the three pairs of Table~\ref{tab:confirm} at 8 blocks of width 256 that train both arms equally long (seeds 20--27
averaged per token). The second column is the share of evaluation tokens. Brackets are 95\%
bootstrap intervals over the 32 evaluation windows (10,000 resamples). Exploratory; the claims are
in Section~\ref{sec:types}.}
\label{tab:types}
\end{table}

\begin{table}[!htb]
\centering
\footnotesize
\setlength{\tabcolsep}{4pt}
% Generated by scripts/paper_results.py from runs/. Do not edit.
\begin{tabular}{@{}lrcccc@{}}
\toprule
Token type & \% & $P=0$ & $P=1$ & $P=2$ & $P=4$ \\
\midrule
All tokens & 100.0 & $+0.280$ & $+0.015$ & $+0.003$ & $+0.001$ \\
 &  & {\scriptsize [$+0.215$, $+0.344$]} & {\scriptsize [$+0.010$, $+0.020$]} & {\scriptsize [$+0.001$, $+0.005$]} & {\scriptsize [$-0.001$, $+0.003$]} \\
\midrule
Exact repeat, far & 16.4 & $+0.646$ & $+0.000$ & $-0.003$ & $-0.001$ \\
 &  & {\scriptsize [$+0.527$, $+0.773$]} & {\scriptsize [$-0.003$, $+0.003$]} & {\scriptsize [$-0.006$, $+0.000$]} & {\scriptsize [$-0.002$, $+0.001$]} \\
Exact repeat, near & 45.2 & $+0.169$ & $+0.001$ & $-0.001$ & $-0.001$ \\
 &  & {\scriptsize [$+0.105$, $+0.255$]} & {\scriptsize [$-0.001$, $+0.002$]} & {\scriptsize [$-0.002$, $+0.001$]} & {\scriptsize [$-0.004$, $+0.000$]} \\
Consistent cue & 12.1 & $+0.465$ & $+0.002$ & $-0.001$ & $-0.001$ \\
 &  & {\scriptsize [$+0.412$, $+0.521$]} & {\scriptsize [$-0.002$, $+0.005$]} & {\scriptsize [$-0.004$, $+0.002$]} & {\scriptsize [$-0.004$, $+0.001$]} \\
Latest value & 2.3 & $+0.262$ & $+0.006$ & $-0.004$ & $-0.005$ \\
 &  & {\scriptsize [$+0.191$, $+0.349$]} & {\scriptsize [$-0.006$, $+0.018$]} & {\scriptsize [$-0.011$, $+0.003$]} & {\scriptsize [$-0.015$, $+0.005$]} \\
Older value & 2.1 & $+0.261$ & $+0.053$ & $+0.012$ & $+0.016$ \\
 &  & {\scriptsize [$+0.209$, $+0.314$]} & {\scriptsize [$+0.036$, $+0.073$]} & {\scriptsize [$-0.004$, $+0.025$]} & {\scriptsize [$+0.003$, $+0.029$]} \\
Novel token & 5.4 & $+0.016$ & $-0.009$ & $+0.008$ & $-0.007$ \\
 &  & {\scriptsize [$-0.001$, $+0.032$]} & {\scriptsize [$-0.021$, $+0.003$]} & {\scriptsize [$-0.001$, $+0.019$]} & {\scriptsize [$-0.017$, $+0.002$]} \\
No match & 5.4 & $+0.048$ & $+0.053$ & $+0.019$ & $+0.013$ \\
 &  & {\scriptsize [$+0.035$, $+0.061$]} & {\scriptsize [$+0.026$, $+0.099$]} & {\scriptsize [$+0.004$, $+0.040$]} & {\scriptsize [$+0.005$, $+0.021$]} \\
Near-duplicate line, far & 3.7 & $+0.261$ & $+0.011$ & $+0.001$ & $+0.003$ \\
 &  & {\scriptsize [$+0.193$, $+0.335$]} & {\scriptsize [$+0.003$, $+0.020$]} & {\scriptsize [$-0.007$, $+0.009$]} & {\scriptsize [$-0.001$, $+0.008$]} \\
Closing bracket, far & 0.1 & $+0.628$ & $+0.074$ & $-0.034$ & $+0.117$ \\
 &  & {\scriptsize [$+0.327$, $+0.963$]} & {\scriptsize [$-0.088$, $+0.273$]} & {\scriptsize [$-0.123$, $+0.055$]} & {\scriptsize [$-0.048$, $+0.288$]} \\
Next integer & 1.0 & $+0.048$ & $+0.004$ & $-0.002$ & $-0.008$ \\
 &  & {\scriptsize [$+0.012$, $+0.091$]} & {\scriptsize [$-0.007$, $+0.016$]} & {\scriptsize [$-0.013$, $+0.008$]} & {\scriptsize [$-0.014$, $-0.002$]} \\
Format & 4.7 & $+0.055$ & $+0.041$ & $+0.012$ & $+0.005$ \\
 &  & {\scriptsize [$+0.041$, $+0.070$]} & {\scriptsize [$+0.013$, $+0.102$]} & {\scriptsize [$+0.000$, $+0.036$]} & {\scriptsize [$-0.002$, $+0.013$]} \\
\bottomrule
\end{tabular}

\caption{Cost of removing both attention layers, $C(P) = L(A{=}0, P) - L(A{=}2, P)$ in nats per
token, by token type and number of placements $P$ (Figure~\ref{fig:types} plots $P = 0$ and $P =
1$). The second column is the share of evaluation tokens of each type. Brackets are 95\% bootstrap
intervals over the 32 evaluation windows (10,000 resamples, windows resampled jointly for both
models, seeds 0 and 1 averaged). Exploratory; the claims are in Section~\ref{sec:types}.}
\label{tab:a3}
\end{table}

\begin{table}[!htb]
\centering
\small
% Generated by scripts/paper_results.py from runs/. Do not edit.
\begin{tabular}{@{}llcccc@{}}
\toprule
& & \multicolumn{3}{c}{Induction share of far attention} & Far attention \\
\cmidrule(lr){3-5}
& & seed 0 & seed 1 & mean & mass (mean) \\
\midrule
$A=2$ & $P=0$ & $0.235$ & $0.231$ & $0.233$ & $0.588$ \\
$A=2$ & $P=1$ & $0.135$ & $0.136$ & $0.136$ & $0.483$ \\
$A=2$ & $P=2$ & $0.157$ & $0.206$ & $0.181$ & $0.471$ \\
$A=2$ & $P=4$ & $0.238$ & $0.110$ & $0.174$ & $0.530$ \\
\midrule
$A=4$ & $P=0$ & $0.245$ & $0.178$ & $0.211$ & $0.500$ \\
$A=4$ & $P=1$ & $0.104$ & $0.205$ & $0.154$ & $0.445$ \\
$A=4$ & $P=2$ & $0.160$ & $0.139$ & $0.150$ & $0.392$ \\
$A=4$ & $P=4$ & $0.072$ & $0.146$ & $0.109$ & $0.420$ \\
\bottomrule
\end{tabular}

\caption{Induction share of far attention on the dose grid, as defined in
Appendix~\ref{app:grid}, and the mean attention mass on keys more than 512 tokens back. Seeds 0
and 1; exploratory.}
\label{tab:a4}
\end{table}

\section{Per-seed losses}
\label{app:seeds}

No run was dropped from any comparison; a run that failed on the platform was relaunched with the
same seed. Tables~\ref{tab:seeds164} and~\ref{tab:seedslong} list every confirmatory run and
Table~\ref{tab:gridseeds} both seeds of every cell of the dose grid.

The first column of Table~\ref{tab:seeds164} is the model with cue lengths 1 to 5. Before we
adopted the longer cues, a pre-registered single test had compared it with the 3:1 hybrid on the
same seeds: $\Delta = -0.117$ (95\% CI $[-0.143, -0.099]$; 8 of 8 seeds; $p = 0.0039$).

\begin{table}[!htb]
\centering
\footnotesize
% Generated by scripts/paper_results.py from runs/. Do not edit.
\begin{tabular}{@{}lcccccc@{}}
\toprule
& \multicolumn{2}{c}{GDN + \ingram{}, $8\times256$} & 3:1 hybrid & 3:1 hybrid & 3:1 + static & GDN + \ingram{} \\
Seed & $N \le 5$ & $N \le 32$ & $8\times256$ & $12\times384$ & table, $8\times256$ & $12\times384$ \\
\midrule
20 & 0.6348 & 0.6129 & 0.8350 & 0.7238 & 0.7618 & 0.5850 \\
21 & 0.6296 & 0.6048 & 0.7409 & 0.7403 & 0.8702 & 0.5829 \\
22 & 0.6307 & 0.6079 & 0.7484 & 0.7093 & 0.7112 & 0.5836 \\
23 & 0.6360 & 0.6009 & 0.7216 & 0.7022 & 0.7458 & 0.5837 \\
24 & 0.6327 & 0.6036 & 0.7237 & 0.6972 & 0.7295 & 0.5818 \\
25 & 0.6287 & 0.6088 & 0.7404 & 0.6919 & 0.7435 & 0.5809 \\
26 & 0.6326 & 0.6101 & 0.7371 & 0.6973 & 0.6729 & 0.5764 \\
27 & 0.6312 & 0.6089 & 0.7473 & 0.7108 & 0.7216 & 0.5825 \\
\midrule
Mean & 0.6320 & 0.6073 & 0.7493 & 0.7091 & 0.7446 & 0.5821 \\
With cache & 0.5942 & 0.5804 & 0.6267 & 0.6029 & 0.6227 & 0.5587 \\
\bottomrule
\end{tabular}

\caption{Loss per seed at 164M training tokens (10,000 steps), for the comparisons of
Table~\ref{tab:confirm} and Appendix~\ref{sec:cues}. The last row is the mean over seeds after
mixing each model with the n-gram cache.}
\label{tab:seeds164}
\end{table}

\begin{table}[!htb]
\centering
\footnotesize
% Generated by scripts/paper_results.py from runs/. Do not edit.
\begin{tabular}{@{}lccccccc@{}}
\toprule
& \multicolumn{2}{c}{328M tokens, $8\times256$} & \multicolumn{2}{c}{328M tokens, $16\times512$} & \multicolumn{3}{c}{1.31B tokens, $8\times256$} \\
\cmidrule(lr){2-3}\cmidrule(lr){4-5}\cmidrule(lr){6-8}
Seed & GDN + \ingram{} & 3:1 & GDN + \ingram{} & 3:1 & 7:1 + \ingram{} & 3:1 & 7:1 \\
\midrule
20 & 0.5505 & 0.6097 & 0.5040 & 0.5386 & 0.4478 & 0.5220 & 0.4976 \\
21 & 0.5496 & 0.6225 & 0.5013 & 0.5572 & 0.4461 & 0.5263 & 0.5090 \\
22 & 0.5517 & 0.5941 & 0.5030 & 0.5533 & 0.4482 & 0.4935 & 0.5298 \\
23 & 0.5480 & 0.6414 & 0.5050 & 0.5943 & 0.4482 & 0.5250 & 0.5113 \\
24 & 0.5482 & 0.6217 &  &  & 0.4521 & 0.5119 & 0.5053 \\
25 & 0.5458 & 0.6490 &  &  & 0.4460 & 0.4995 & 0.5082 \\
26 & 0.5535 & 0.6282 &  &  & 0.4442 & 0.5092 & 0.4965 \\
27 & 0.5513 & 0.6121 &  &  & 0.4480 & 0.4950 & 0.5053 \\
\midrule
Mean & 0.5498 & 0.6223 & 0.5033 & 0.5609 & 0.4476 & 0.5103 & 0.5079 \\
With cache & 0.5302 & 0.5544 & 0.4879 & 0.5090 & 0.4375 & 0.4738 & 0.4718 \\
\bottomrule
\end{tabular}

\caption{Loss per seed for the longer runs of Table~\ref{tab:confirm}.
The third size ran on seeds 20--23 only.}
\label{tab:seedslong}
\end{table}

\begin{table}[!htb]
\centering
\small
% Generated by scripts/paper_results.py from runs/. Do not edit.
\begin{tabular}{@{}lcccc@{}}
\toprule
Attention layers & $P=0$ & $P=1$ & $P=2$ & $P=4$ \\
\midrule
$A=0$ (GDN) & $1.002$ & $0.609$ & $0.590$ & $0.584$ \\
$A=1$ (7:1) & $0.995$ & $0.607$ & $0.585$ & $0.582$ \\
$A=2$ (3:1) & $0.722$ & $0.594$ & $0.587$ & $0.583$ \\
$A=4$ (1:1) & $0.742$ & $0.602$ & $0.584$ & $0.575$ \\
\bottomrule
\end{tabular}

\caption{Seed-mean loss on the dose grid of Figure~\ref{fig:grid} (seeds 0 and 1). Cells with
$P = 0$ train at learning rate $4 \times 10^{-3}$ and cells with $P \ge 1$ at $2 \times 10^{-3}$
(Table~\ref{tab:lr}).}
\label{tab:grid}
\end{table}

\begin{table}[!htb]
\centering
\footnotesize
% Generated by scripts/paper_results.py from runs/. Do not edit.
\begin{tabular}{@{}lcccc@{}}
\toprule
Attention layers & $P=0$ & $P=1$ & $P=2$ & $P=4$ \\
\midrule
$A=0$ & 1.0175 / 0.9861 & 0.6109 / 0.6064 & 0.5914 / 0.5894 & 0.5850 / 0.5830 \\
$A=1$ & 0.9965 / 0.9940 & 0.6068 / 0.6068 & 0.5901 / 0.5808 & 0.5806 / 0.5836 \\
$A=2$ & 0.7141 / 0.7295 & 0.5974 / 0.5903 & 0.5883 / 0.5865 & 0.5846 / 0.5812 \\
$A=4$ & 0.7568 / 0.7268 & 0.6054 / 0.5982 & 0.5832 / 0.5842 & 0.5771 / 0.5732 \\
\bottomrule
\end{tabular}

\caption{Loss on the dose grid, seed 0 / seed 1 (means in Table~\ref{tab:grid}).}
\label{tab:gridseeds}
\end{table}

\begin{table}[!htb]
\centering
\small
% Generated by scripts/paper_results.py from runs/. Do not edit.
\begin{tabular}{@{}lccccc@{}}
\toprule
Cue lengths $N$ & Parameters & Seed 0 & Seed 1 & Mean & Mean with cache \\
\midrule
1, 2, 3, 4, 5 & 51.7M & 0.6349 & 0.6316 & $0.633$ & $0.595$ \\
1, 2, 3, 4, 5, 8 & 56.0M & 0.6237 & 0.6217 & $0.623$ & $0.592$ \\
1, 2, 3, 5, 8 & 51.7M & 0.6243 & 0.6244 & $0.624$ & $0.593$ \\
1, 2, 4, 8, 16 & 51.7M & 0.6161 & 0.6124 & $0.614$ & $0.586$ \\
1, 2, 3, 5, 8, 16 & 56.0M & 0.6111 & 0.6112 & $0.611$ & $0.583$ \\
1, 2, 4, 8, 16, 32 & 56.0M & 0.6102 & 0.6020 & $0.606$ & $0.579$ \\
\bottomrule
\end{tabular}

\caption{Loss by cue set for GDN with \ingram{} ($P = 1$, 8 blocks of width 256, 164M
tokens), seeds 0 and 1 (exploratory; used to select the cue set, Appendix~\ref{sec:cues}). Each cue
length adds one table of $2^{16}$ rows of width 64. The last column interpolates the model with
the n-gram cache.}
\label{tab:cues}
\end{table}

\section{Learning-rate checks}
\label{app:lr}

For each model of Table~\ref{tab:lr} we compared two or three learning rates on the selection
seeds, and each confirmatory arm in the table uses the rate with the lowest mean. The 7:1 hybrid
and the 7:1 hybrid with \ingram{} have no comparison of their own: they use
$4 \times 10^{-3}$, the rate of the 3:1 hybrid, and $2 \times 10^{-3}$, the rate of GDN with
\ingram{}, as do the remaining cells of the dose grid without and with \ingram{}.
\ingram{}'s rate of $2 \times 10^{-3}$ was tuned with cue lengths 1 to 5 and fixed before the runs
with long cues; the rows for long cues are a sensitivity check and selected nothing. At 16 blocks
of width 512 each rate has one seed. The 80,000-step runs of Section~\ref{sec:length} reuse the
rates selected at 10,000 steps.

\begin{table}[!htb]
\centering
\footnotesize
% Generated by scripts/paper_results.py from runs/. Do not edit.
\begin{tabular}{@{}llccc@{}}
\toprule
Model & Learning rate & Seed 0 & Seed 1 & Mean \\
\midrule
GDN, $8\times256$ & 4e-3 & 1.0175 & 0.9861 & $1.002$ \\
 & 8e-3 & 1.1702 & 1.0317 & $1.101$ \\
3:1 hybrid, $8\times256$ & 2e-3 & 0.7329 & 0.7175 & $0.725$ \\
 & 4e-3 & 0.7141 & 0.7295 & $0.722$ \\
1:1 hybrid, $8\times256$ & 2e-3 & 0.7439 & 0.7202 & $0.732$ \\
 & 4e-3 & 0.7568 & 0.7268 & $0.742$ \\
Attention only, $8\times256$ & 2e-3 & 0.7905 & 0.7408 & $0.766$ \\
 & 4e-3 & 0.8148 & 0.7847 & $0.800$ \\
GDN + \ingram{}, $8\times256$ & 1e-3 & 0.6239 & 0.6218 & $0.623$ \\
 & 2e-3 & 0.6102 & 0.6020 & $0.606$ \\
 & 4e-3 & n/a & 0.6106 & n/a \\
3:1 hybrid + static table, $8\times256$ & 2e-3 & 0.7293 & 0.7554 & $0.742$ \\
 & 4e-3 & 0.7425 & 0.7333 & $0.738$ \\
3:1 hybrid, $12\times384$ & 2e-3 & 0.6772 & 0.7481 & $0.713$ \\
 & 4e-3 & 0.7139 & 0.7767 & $0.745$ \\
GDN + \ingram{} (cues 1--5), $12\times384$ & 1e-3 & 0.6127 & 0.6175 & $0.615$ \\
 & 2e-3 & 0.6078 & 0.6093 & $0.609$ \\
3:1 hybrid, $16\times512$ & 1e-3 & 0.5692 & n/a & n/a \\
 & 2e-3 & 0.5540 & n/a & n/a \\
GDN + \ingram{}, $16\times512$ & 1e-3 & 0.5038 & n/a & n/a \\
 & 2e-3 & 0.5017 & n/a & n/a \\
\bottomrule
\end{tabular}

\caption{Loss by learning rate on the selection seeds, at 164M training tokens (328M for 16 blocks
of width 512), 16k context. An entry of n/a marks a run that was not made or was lost to a host
failure.}
\label{tab:lr}
\end{table}

\section{Read rule, quantization and generation}
\label{app:other}

All results above use the additive read; in a smaller setting (4 blocks of width 256, 4,096-token
context, the first 4,096 trajectories, 4 windows per step, 128 evaluation windows, cue lengths 1
to 5, seeds 13--20), we pre-registered a single test of a scored read,
which stores $K$ notes per slot and weights them by a key computed from the context. With $K = 4$
it lowers the loss from 0.829 to 0.824, $\Delta = -0.005$ ($[-0.008, -0.003]$; 8 of 8 seeds;
$p = 0.0039$), and stores 8 vectors of state per slot where the additive read stores one. The
difference with $K = 8$ is $-0.012$ ($[-0.015, -0.009]$; 8 of 8; reported, not tested). On the
recipe of this paper (long cues, 16k context), the scored read with $K = 8$ was worse than the
additive read on both selection seeds (0.629 and 0.608 against 0.610 and 0.602; one learning rate).

We also checked quantization on checkpoints from seeds 0 and 1 that use cue lengths 1 to 5
(exploratory). Under round-to-nearest quantization applied at evaluation, 4-bit rows of the hashed
table change the loss of GDN with \ingram{} by $+0.0001$, and rounding the slot state to 4 bits at
each read by $+0.0005$; the writes are computed at full precision. With 4-bit weights in the
linear layers whose input width is a multiple of 64, and 8-bit weights in the embedding and the
output layer, that model loses $+0.003$ and $+0.002$ on the two seeds, against $+0.007$ and
$+0.008$ for GDN and $+0.010$ and $+0.009$ for the 3:1 hybrid; the read projections $W_K$ and $W_O$ of
\ingram{}, of input width 325 in these checkpoints, remain at full precision. Keys and values of the hybrid at 8 bits cost at most $+0.001$. Our
4-bit keys and values are quantized per token, without per-channel scaling, and cost $+0.350$ and
$+0.064$ on the two seeds.

For generation we provided the same checkpoints with 32 prompts of 4,096 tokens and sampled 256
new tokens.
We classify a continuation as a loop when its last 128 tokens are periodic with a period of at most
64. The reference continuations loop in 7 of the 32 prompts. At temperature 1, every model loops
in 7 to 9. Under greedy decoding the counts are 16 to 21 for the models with \ingram{}, 16 and 19
for the 3:1 hybrid, and 18 and 20 for GDN. We detect no increase in looping with \ingram{}
at 32 prompts and two seeds. Greedy continuations of the models with \ingram{} share more 8-grams
with the prompt (0.69 to 0.70) than those of the hybrid (0.58 and 0.64) or the references (0.62).

\section{Decode benchmark}
\label{app:decode}

We timed token-by-token decoding once, on an RTX~4090 at batch size 1, for untrained models of 4
blocks and width 256 on tokens of WikiText-103: 256 decoded tokens after a prefill of 4,096 to
1,048,576 tokens. The \ingram{} model in this benchmark is an earlier configuration with ten slot
streams (cue lengths 1 to 5 under two write rules) of slot width 32, placed on the embeddings,
with a separate static table of $2^{17}$ rows per stream that has its own gate in place of the
additive table row. Per decoded token after 4,096 and 1,048,576
tokens of context, GDN took 4.91 and 4.90~ms, GDN with \ingram{} 6.22 and 6.26~ms, a hybrid with
one attention layer 4.78 and 5.79~ms, and a model of four attention layers 4.06 and 5.97~ms.
\ingram{} thus adds 1.3 to 1.4~ms per token at both lengths; the benchmark records one total per
model and does not separate the host-side hashing from the reads and writes on the accelerator.

\section{Compute}
\label{app:compute}

Every training run reported in this paper used one rented NVIDIA RTX~4090 with 24~GB of memory.
The machines came from the Vast.ai marketplace and were rented through Rhombic
(\url{https://rhombic.space}), an orchestration service that launched the runs, tracked their cost
and collected their logs. On hosts of different speed, a run of 10,000 steps took 17 to 33
minutes, a run of 20,000 steps 39 to 74 minutes, and a run of 80,000 steps 2.1 to 3.6 hours. The
88 runs behind Figure~\ref{fig:confirm} took 103 GPU-hours of training and were billed at 64 US
dollars with their relaunches, and a run of 80,000 steps cost 1.49 dollars on average. All 534
logged runs of the project, which include the selection runs and exploratory runs that this paper
does not report, took 319 hours of training time. The account that rented the machines was billed
213 dollars in total, for 700 rentals and 420 hours at 0.51 dollars per hour on average, which is
an upper bound on the cost of the project. Data preparation, the token statistics of Section~\ref{sec:repeats} and all analyses
ran on CPUs.

\section{Use of AI tools}
\label{app:ai}

The author proposed the method, the research questions and the framing of this paper, decided which
experiments to run, reviewed the results and the text, and is responsible for the content of the
paper. An AI coding assistant (Claude,
Anthropic, used through Claude Code) wrote most of the code under the author's direction, launched
and monitored the training runs, wrote the analysis scripts, searched and summarized the related
work, and drafted the text and the figures. Three practices limit the errors that this division of
work can introduce. Each confirmatory comparison was committed to the repository as a
pre-registration before its runs started. The numbers in the results are produced by scripts from
the run logs and are not typed by hand. The descriptions of the most closely related papers were
checked against the full text of each paper.
